% !TeX program = xelatex
% UTF-8. Standalone XeLaTeX/XePersian, or input into a XePersian document.
% Transcription of Notebook/ped seizure; source medical statements are preserved.
% pic01.jpg = 20261016_210320778.jpg
% pic02.jpg = 20261016_210327004.jpg
% pic03.jpg = 20261016_210331567.jpg
% pic04.jpg = 20261016_210339046.jpg
% pic05.jpg = 20261016_210344156.jpg
% pic06.jpg = 20261016_210350177.jpg
% pic07.jpg = 20261016_210356250.jpg
% pic08.jpg = 20261016_210400605.jpg
\ifdefined\settextfont
  \def\PedSeizureStandalone{0}
\else
  \def\PedSeizureStandalone{1}
  \documentclass[11pt,a4paper]{article}
  \usepackage[margin=2cm]{geometry}
  \usepackage{amsmath,amssymb,enumitem}
  \usepackage{xepersian}
  \IfFontExistsTF{Amiri}{\settextfont{Amiri}}{\settextfont{FreeSerif}}
  \setlatintextfont{FreeSerif}
  \setlist[itemize]{itemsep=2pt,topsep=3pt,parsep=0pt}
\fi
\providecommand{\en}[1]{\lr{#1}}
\newcommand{\psarrow}{\ensuremath{\leftarrow}}
\newcommand{\psunclear}[1]{\textbf{[خوانش نامطمئن: #1]}}
\newcommand{\pspage}[2]{\clearpage\section*{۱۸۱ ـ تشنج — صفحهٔ #1}
  {\footnotesize\noindent\en{#2}\par
  \noindent\en{Date:}\hfill\en{Subject:} ۱۸۱ ـ تشنج\par}
  \medskip}
\newcommand{\psfooter}[1]{\par\smallskip{\footnotesize\en{GNOTE}\hfill #1\par}}
\ifnum\PedSeizureStandalone=1 \begin{document}\fi

\pspage{۱}{pic01.jpg — 20261016\_210320778.jpg}
\noindent کودکان
\subsection*{اپیلپسی}
\begin{itemize}
\item \en{Focal}: از یک بخش خاص از کورتکس منشأ می‌گیرد.
\item \en{generalized}: از هر دو نیمکره.
\end{itemize}
\begin{latin}
\noindent Focal seizures $\rightarrow$ Motor abnormalities $\rightarrow$ Tonic, clonic, Myoclonic\par
\noindent Focal seizures $\rightarrow$ Sensory\par
\noindent Focal seizures $\rightarrow$ Psychic\par
\noindent Focal seizures $\rightarrow$ Autonomic
\end{latin}
\subsection*{فوکال ساده با حفظ هوشیاری (\en{Simple partial})}
\begin{itemize}
\item از یک منطقهٔ آناتومیک خاص؛ می‌تواند به نواحی اطراف گسترش یابد یا نیابد.
\end{itemize}
\subsection*{تشنج کمپلکس با اختلال هوشیاری (\en{Complex partial})}
\begin{itemize}
\item تشنج کمپلکس به همراه اختلال سطح هوشیاری است.
\item هوشیاری ممکن است کاملاً از دست نرود.
\item ممکن است \en{Stare}، \en{Automatisms} \psarrow{} حرکات غیرهدفمند دهان و اندام‌ها.
\item درمان: لوتیراستام، اکس‌کاربازپین، لاموتریژین.
\item تمپورال لوب سیزر \psarrow{} دژاوو (\en{psychic symptomes}).
\item زمانی که یک \en{focal seizure} به هر دو نیمکره گسترش یابد تبدیل به
\en{Bilateral tonic-clonic seizure (secondarily generalized seizure)}.
\end{itemize}
\subsection*{تشنج جنرالیزه}
\begin{itemize}
\item تشنج‌های تونیک، کلونیک یا تونیک ـ کلونیک به تنهایی یا با انواع دیگر تشنج.
\item به طور معمول تشنج به صورت ناگهانی شروع می‌شود؛ ولی می‌تواند به دنبال یک سری \en{Myoclonic jerks} باشد.
\item در یک تشنج تونیک ـ کلونیک، هوشیاری و کنترل عضلات از دست می‌رود.
\end{itemize}
\begin{latin}
\noindent Tonic stiffening + Upward deviation of the eyes
\end{latin}
\noindent \en{HTN}، تجمع ترشحات، میدریاز، دیافورز.
\psfooter{۱}

\pspage{۲}{pic02.jpg — 20261016\_210327004.jpg}
\begin{itemize}
\item \en{Clonic jerks} به دنبال فاز تونیک رخ می‌دهد.
\item فاز پست‌ایکتال کودک ممکن است هایپوتون باشد.
\item پس از بیدار شدن \psarrow{} سردرد و تحریک‌پذیری.
\end{itemize}
\subsection*{تشنج \en{Absence}}
\begin{latin}
\noindent Hallmark: Brief ($<10$ s) loss of awareness + Eyelid fluttering\par
\noindent or simple automatisms
\end{latin}
\begin{itemize}
\item معمولاً ظرف ۴–۶ سالگی شروع می‌شود.
\item تصویربرداری و ارزیابی‌ها \en{N/L} است.
\end{itemize}
\begin{latin}
\noindent EEG: generalized 3-Hz spike-and-wave activity
\end{latin}
\begin{itemize}
\item هایپرونتیلیشن آن را تحریک می‌کند.
\item فاز پست‌ایکتال ندارد \psarrow{} چندین تشنج در روز.
\item خط اول درمان: والپروئیک اسید، اتوسوکسیمید.
\item صرع مایوکلونیک معمولاً با چندین نوع تشنج همراهی دارد.
\end{itemize}
\subsection*{\en{Febrile seizure}}
\begin{itemize}
\item ممکن است ناشی از عفونت \en{CNS} باشد (انسفالیت، مننژیت و آبسهٔ مغزی).
\item شایع‌ترین علت تشنج در کودکان ۶ ماه تا ۵ سالگی است.
\item \en{Simple febrile seizure} \psarrow{} پروگنوز عالی.
\item جنرالیزه، زیر ۱۵ دقیقه، یک بار در ۲۴ ساعت.
\item \en{Complex febrile seizure}: ویژگی‌های فوکال، بالای ۱۵ دقیقه، چند بار در ۲۴ ساعت.
\end{itemize}
\psfooter{۲}

\pspage{۳}{pic03.jpg — 20261016\_210331567.jpg}
\subsection*{\en{Febrile seizure} (ادامه)}
\begin{itemize}
\item در ۳۰–۵۰\% عود می‌کند ولی به هوش آسیب نمی‌زند.
\item ریسک تبدیل به صرع به اندازهٔ جامعهٔ عادی است.
\item موارد \en{Simple}: در طول تشنج \psarrow{} دیازپام رکتال.
\item درمان دیگری لازم نیست.
\item داروهای تب‌بر کمک‌کننده است.
\end{itemize}
\subsection*{\en{Psychogenic nonepileptic events (PNEE)} (سودوسیژر)}
\begin{itemize}
\item در اختلال کانورژن و سوماتیک سیمپتوم دیس‌اوردر دیده می‌شود.
\item طی حمله چشم‌های بیمار بسته است.
\item حرکات بدن: \en{thrashing}، \en{tremulousness}.
\item \en{Verbalization}، \en{Pelvic thrusting}.
\item بی‌اختیاری ادرار و مدفوع پیدا نمی‌کند.
\item نوک زبان را گاز می‌گیرد.
\item \en{EEG} طی حمله \en{N/L} است.
\end{itemize}
\subsection*{سندرم‌های اپیلپسی}
\begin{itemize}
\item اگر کودک مبتلا به \en{ADHD} و صرع داشته باشه می‌توانیم با محرک‌ها درمان کنیم.
\end{itemize}
\subsection*{۱ ـ \en{Benign Rolandic Epilepsy}}
\begin{latin}
\noindent (childhood epilepsy with centrotemporal spikes)
\end{latin}
\begin{itemize}
\item شروع بین ۵–۱۰ سالگی.
\item در زمان خواب یا حین بیدار شدن.
\item تشنج‌های فوکال موتور: \en{face}، \en{arm}.
\item \en{drooling}، \en{rhythmic guttural sound}، اختلال بلع، \en{speech}.
\end{itemize}
\begin{latin}
\noindent EEG: N/L or independent bilateral centrotemporal sharp waves
\end{latin}
\begin{itemize}
\item برای تشخیص، تصویربرداری لازم نیست.
\item معاینهٔ نورولوژیک \en{N/L} است؛ به خوبی به درمان جواب می‌دهد.
\item هوش \en{N/L} است.
\item تشنج، صرع پس از بلوغ از بین می‌رود.
\end{itemize}
\psfooter{۳}

\pspage{۴}{pic04.jpg — 20261016\_210339046.jpg}
\subsection*{۲ ـ \en{Childhood absence epilepsy}}
\begin{itemize}
\item در سنین مدرسه شروع می‌شود \psarrow{} در نوجوانی قطع می‌شود.
\item اگر قطع نشود در برخی از موارد به \en{JME} تبدیل می‌شود.
\item اتوسوکسیمید خط اول درمان است.
\item برخی از بیماران تشنج تونیک ـ کلونیک جنرالیزه هم دارند \psarrow{} والپروئیک اسید.
\end{itemize}
\subsection*{۳ ـ \en{JME}}
\begin{itemize}
\item شایع‌ترین صرع جنرالیزه در نوجوانان و جوانان.
\item شروع معمولاً در نوجوانی و با \en{myoclonic jerks}.
\item صبح‌ها \psarrow{} وسایل از دست می‌افتند.
\item حاشیهٔ صفحه: \psunclear{۱۷۵}.
\item به خوبی به درمان ضدصرع جواب می‌دهد.
\item خط اول والپروئیک اسید (تراتـوژن است) \psarrow{} لوتیراستام (برای دختران).
\item تا آخر عمر باید درمان بگیرند.
\end{itemize}
\subsection*{۴ ـ \en{Infantile spasms}}
\begin{itemize}
\item ۳–۸ ماهگی شایع‌ترین سن؛ یک اورژانس نورولوژیک است.
\end{itemize}
\begin{latin}
\noindent Brief contractions of the neck, trunk, arm\par
\noindent $\rightarrow$ followed by a phase of sustained muscle contraction\par
\noindent $\rightarrow$ less than 2 seconds
\end{latin}
\begin{itemize}
\item زمانی که شیرخوار از خواب بیدار می‌شود یا می‌خواهد بخوابد \psarrow{} شایع‌ترین زمان.
\item حاشیهٔ صفحه: \psunclear{علامت کوتاه شبیه «۹+»}.
\end{itemize}
\begin{latin}
\noindent Each jerk is followed by a brief period of relaxation
\end{latin}
\subsection*{۵ ـ سندرم \en{West}}
\begin{itemize}
\item تریاد: \en{infantile spasms}، \en{developmental regression}، \en{EEG} به شدت غیرطبیعی.
\end{itemize}
\begin{latin}
\noindent Hypsarrhythmia (chaotic high-voltage slow waves, spikes and polyspikes)
\end{latin}
\begin{itemize}
\item اسپاسم شیرخوارگی بیش از ۲۰۰ اتیولوژی دارد.
\item تروما، توبروس اسکلروزیس، استروک، فنیل‌کتون‌اوری، \ldots
\end{itemize}
\psfooter{۴}

\pspage{۵}{pic05.jpg — 20261016\_210344156.jpg}
\begin{itemize}
\item شیرخوارانی که اتیولوژی اسپاسم آن‌ها شناسایی می‌شود، ریسک بسیار بالایی برای مشکلات نورولوژیک در آینده دارند.
\item درمان اسپاسم شیرخوارگی: آدرنوکورتیکوتروپیک هورمون، \en{OCS} دوز بالا، ویگاباترین.
\item اگر علت زمینه‌ای توبروس اسکلروزیس باشد \psarrow{} ویگاباترین.
\end{itemize}
\subsection*{۶ ـ \en{Lennox--Gastaut Syndrome}}
\begin{itemize}
\item یک سندرم صرعی شدید است؛ سن شروع متفاوت ولی اغلب زیر ۵ سال.
\end{itemize}
\begin{latin}
\noindent Multiple seizure types: Atonic, focal, atypical, generalized
\end{latin}
\begin{itemize}
\item اکثر بیماران اختلال زمینه‌ای مغزی دارند یا ژنتیکی.
\item تشنج‌ها به سختی کنترل می‌شوند.
\item بیماران معمولاً ناتوانی ذهنی شدید دارند.
\end{itemize}
\subsection*{۷ ـ \en{Benign (self-limited) neonatal convulsions}}
\begin{itemize}
\item وراثت \en{AD}؛ اختلال در کانال پتاسیم نورون.
\item نوزاد سالم با تشنج فوکال که تا هفتهٔ اول تولد طول می‌کشد؛ \en{fifth-day fits}.
\item پاسخ به درمان خوب است و پروگنوز مطلوب می‌باشد.
\end{itemize}
\subsection*{۸ ـ \en{Acquired epileptic aphasia}}
\begin{latin}
\noindent (Landau--Kleffner syndrome)
\end{latin}
\begin{itemize}
\item از دست رفتن ناگهانی توانایی‌های زبانی.
\item یک آگنوزی شنوایی، یک نقص شنوایی کورتیکال.
\item نادر در پسران.
\end{itemize}
\begin{latin}
\noindent EEG: Highly epileptiform in sleep
\end{latin}
\begin{itemize}
\item کیس: بیمار جوان با رگرسیون زبانی واضح.
\item \psunclear{به کلی} درمان‌پذیر است.
\end{itemize}
\psfooter{۵}

\pspage{۶}{pic06.jpg — 20261016\_210350177.jpg}
\subsection*{استاتوس اپی‌لپتیکوس}
\begin{itemize}
\item یک اورژانس نورولوژیک است.
\item الف) تشنج ادامه‌دار برای حداقل ۵ دقیقه؛ ب) تشنج‌های مکرر بدون فاز پست‌ایکتال.
\item ۱۴\% ریسک عوارض نورولوژیک، ۵\% ریسک مرگ.
\item اتیولوژی: مسمومیت دارویی، هایپوکلسمی و اختلال الکترولیت، تروما به سر، قطع داروی ضدصرع (فراموش شدن دوز ضدصرع در بیمار صرع)، خونریزی مغزی، استروک، هایپوکسی، اختلال متابولیک.
\item برخورد با بیمار: \en{ABC}، \en{V/S}، \en{O$_2$ therapy}.
\item درمان اولیه: بنزودیازپین؛ وقتی تشنج بیش از ۵ دقیقه طول کشد.
\item دوز مناسب با فاصلهٔ یک بار تکرار.
\item اگر با ۲ دوز بنزودیازپین کنترل نشد:
\item نوزاد و شیرخوار: فنوباربیتال.
\item کودکان بزرگ‌تر: فنی‌توئین، والپروئیک اسید، لوتیراستام.
\item خط سوم: میدازولام و پنتوباربیتال.
\item خط چهارم: بیهوشی عمومی.
\end{itemize}
\subsection*{بررسی}
\begin{itemize}
\item بیمار تب و تشنج که کاملاً ریکاوری یافته نیاز به بررسی بیشتری ندارد.
\item یک کودک سالم با یک نوبت تشنج که خودبه‌خود رفع شده و معاینهٔ نورولوژیک طبیعی است \psarrow{} بررسی دیگری لازم نیست.
\item اگر شک به مننژیت \psarrow{} \en{LP} می‌گیریم.
\item تصویربرداری انتخابی \en{MRI} هست.
\item در بیمار \en{JME} و ابسنس، \en{MRI} نیاز نیست.
\end{itemize}
\psfooter{۶}

\pspage{۷}{pic07.jpg — 20261016\_210356250.jpg}
\begin{itemize}
\item تشنج ابسنس، اسپاسم شیرخواران، آتپیکال ابسنس، تشنج‌های آتونیک \psarrow{} خطر تکرار تشنج بالاست.
\item در اولین نوبت، درمان را آغاز می‌کنیم.
\item در اکثر کودکان: پس از ۲ سال بدون علامت بودن می‌توان دارو را قطع کرد.
\item به جز \en{JME}، صرع میوکلونیک پیشرونده، آتپیکال ابسنس، سندرم لنوکس ـ گاستات.
\item این کودکان می‌توانند در ورزش‌های \psunclear{بدنی/رزمی} شرکت کنند.
\end{itemize}
\subsection*{داروها}
\begin{itemize}
\item صرفاً در تشنج جنرالیزه و ابسنس: اتوسوکسیمید.
\item صرفاً در تشنج فوکال: اکس‌کاربازپین، کاربامازپین، فنی‌توئین.
\item مشترک: لاموتریژین، فنوباربیتال، توپیرامات، والپروئیک اسید.
\end{itemize}
\subsection*{درمان استاتوس (\en{SE})}
\noindent دقیقهٔ ۰–۵: پایدارسازی.
\begin{latin}
\noindent ABC, V/S, O$_2$, BS, EKG, IV\par
\medskip
\noindent if BS $<60$ $\rightarrow$ Adults: 100 mg B$_1$ IV + 50 ml D50W IV\par
\noindent child $<2$: 4 ml/kg D12.5W IV\par
\noindent $2\leq$ child: 2 ml/kg D25W IV
\end{latin}
\noindent دقیقهٔ ۵–۲۰: فاز درمان اولیه.
\begin{latin}
\begin{itemize}
\item Midazolam IM (10 mg for $>40$ kg, 5 mg for 13--40) Single dose
\item Lorazepam IV (0.1 mg/kg/dose, max 4 mg) May repeat once
\item Diazepam IV (0.15--0.2 mg/kg/dose, max 10 mg) May repeat once
\end{itemize}
\end{latin}
\noindent اگر هیچ‌کدام در دسترس نبود:
\begin{latin}
\begin{itemize}
\item Diazepam Rectal (0.2--0.5 mg/kg, max 20) Single dose
\item Phenobarb IV (15 mg/kg/dose) Single dose
\item Midazolam intranasal
\end{itemize}
\end{latin}
\psfooter{۷}

\pspage{۸}{pic08.jpg — 20261016\_210400605.jpg}
% Source page 8 has no handwritten subject or page number; only GNOTE is printed.
\subsection*{دقیقهٔ ۲۰–۴۰: فاز درمان ثانویه}
\begin{latin}
\begin{itemize}
\item Na-Val IV (40 mg/kg, max 3000) Single dose
\item Levebel IV (60 mg/kg, max 4500) Single dose
\item Fosphenytoin IV (20 mg PE/kg, max 1500) Single dose
\end{itemize}
\end{latin}
\noindent اگر هیچ‌کدام در دسترس نبود و تا الان فنوبارب نداده بودیم:
\begin{latin}
\noindent Phenobarb IV (15 mg/kg)
\end{latin}
\subsection*{دقیقهٔ ۴۰–۶۰: فاز درمان ثالثیه}
\begin{itemize}
\item تکرار خط دوم یا:
\item دوز بیهوشی از: تیوپنتال، میدازولام، فنوبارب، پروپوفول.
\item دوز متصل با فلش به میدازولام: \en{0.05--2 mg/kg/h}.
\end{itemize}
\subsection*{دوزهای نگهدارنده}
\begin{itemize}
\item فنوبارب: نگهدارنده \en{3--5 mg/kg/day divided BD}.
\item کما: اینداکشن \en{5--15 mg/kg}؛ نگهدارنده \en{1--5 mg/kg/h}.
\item لوبل: نگهدارنده \en{30--60 mg/kg/day divided BD}؛ ماکزیمم \en{3 gr/day}.
\item والپروات: نگهدارنده \en{30--60 mg/kg/day div BD}.
\end{itemize}
\begin{latin}
\noindent Lab data in SE $\rightarrow$ BS, Na, Ca, Mg, CBC; basic metabolic, CT scan, EEG monitoring
\end{latin}
\par\smallskip{\footnotesize\en{GNOTE}\par}

\ifnum\PedSeizureStandalone=1
  \def\PedSeizureFinish{\end{document}}
\else
  \let\PedSeizureFinish\relax
\fi
\PedSeizureFinish
